%Input preamble
%Style
\documentclass[12pt]{article}
\usepackage[top=1in, bottom=1in, left=1in, right=1in]{geometry}
\parindent 22pt
\usepackage{fancyhdr}

%Packages
\usepackage{adjustbox}
\usepackage{amsmath}
\usepackage{amsfonts}
\usepackage{amssymb}
\usepackage{bm}
\usepackage[table]{xcolor}
\usepackage{tabu}
\usepackage{color,soul}
\usepackage{makecell}
\usepackage{longtable}
\usepackage{multirow}
\usepackage[normalem]{ulem}
\usepackage{etoolbox}
\usepackage{graphicx}
\usepackage{tabularx}
\usepackage{ragged2e}
\usepackage{booktabs}
\usepackage{caption}
\usepackage{fixltx2e}
\usepackage[para, flushleft]{threeparttablex}
\usepackage[capposition=top,objectset=centering]{floatrow}
\usepackage{subcaption}
\usepackage{pdfpages}
\usepackage{pdflscape}
\usepackage{natbib}
\usepackage{bibunits}
\definecolor{maroon}{HTML}{990012}
\usepackage[colorlinks=true,linkcolor=maroon,citecolor=maroon,urlcolor=maroon,anchorcolor=maroon]{hyperref}
\usepackage{marvosym}
\usepackage{makeidx}
\usepackage{tikz}
\usetikzlibrary{shapes}
\usepackage{setspace}
\usepackage{enumerate}
\usepackage{rotating}
\usepackage{tocloft}
\usepackage{epstopdf}
\usepackage[titletoc]{appendix}
\usepackage{framed}
\usepackage{comment}
\usepackage{xr}
\usepackage{titlesec}
\usepackage{footnote}
\usepackage{longtable}
\newlength{\tablewidth}
\setlength{\tablewidth}{9.3in}
\setcounter{secnumdepth}{4}

\titleformat{\paragraph}
{\normalfont\normalsize\bfseries}{\theparagraph}{1em}{}
\titlespacing*{\paragraph}
{0pt}{3.25ex plus 1ex minus .2ex}{1.5ex plus .2ex}
\makeatletter
\pretocmd\start@align
{%
  \let\everycr\CT@everycr
  \CT@start
}{}{}
\apptocmd{\endalign}{\CT@end}{}{}
\makeatother
%Watermark
\usepackage[printwatermark]{xwatermark}
\usepackage{lipsum}
\definecolor{lightgray}{RGB}{220,220,220}
%\newwatermark[allpages,color=lightgray,angle=45,scale=3,xpos=0,ypos=0]{Preliminary Draft}

%Further subsection level
\usepackage{titlesec}
\setcounter{secnumdepth}{4}
\titleformat{\paragraph}
{\normalfont\normalsize\bfseries}{\theparagraph}{1em}{}
\titlespacing*{\paragraph}
{0pt}{3.25ex plus 1ex minus .2ex}{1.5ex plus .2ex}

\setcounter{secnumdepth}{5}
\titleformat{\subparagraph}
{\normalfont\normalsize\bfseries}{\thesubparagraph}{1em}{}
\titlespacing*{\subparagraph}
{0pt}{3.25ex plus 1ex minus .2ex}{1.5ex plus .2ex}

%Functions
\DeclareMathOperator{\cov}{Cov}
\DeclareMathOperator{\corr}{Corr}
\DeclareMathOperator{\var}{Var}
\DeclareMathOperator{\plim}{plim}
\DeclareMathOperator*{\argmin}{arg\,min}
\DeclareMathOperator*{\argmax}{arg\,max}

%Math Environments
\newtheorem{theorem}{Theorem}
\newtheorem{claim}{Claim}
\newtheorem{condition}{Condition}
\renewcommand\thecondition{C--\arabic{condition}}
\newtheorem{algorithm}{Algorithm}
\newtheorem{assumption}{Assumption}
\renewcommand\theassumption{A--\arabic{assumption}}
\newtheorem{remark}{Remark}
\renewcommand\theremark{R--\arabic{remark}}
\newtheorem{definition}[theorem]{Definition}
\newtheorem{hypothesis}[theorem]{Hypothesis}
\newtheorem{property}[theorem]{Property}
\newtheorem{example}[theorem]{Example}
\newtheorem{result}[theorem]{Result}
\newenvironment{proof}{\textbf{Proof:}}{$\bullet$}

%Commands
\newcommand\independent{\protect\mathpalette{\protect\independenT}{\perp}}
\def\independenT#1#2{\mathrel{\rlap{$#1#2$}\mkern2mu{#1#2}}}
\newcommand{\overbar}[1]{\mkern 1.5mu\overline{\mkern-1.5mu#1\mkern-1.5mu}\mkern 1.5mu}
\newcommand{\equald}{\ensuremath{\overset{d}{=}}}
\captionsetup[table]{skip=10pt}
%\makeindex

\setlength\parindent{0pt}
\setlength{\parskip}{10pt}

\newcolumntype{L}[1]{>{\raggedright\let\newline\\\arraybackslash\hspace{0pt}}m{#1}}
\newcolumntype{C}[1]{>{\centering\let\newline\\\arraybackslash\hspace{0pt}}m{#1}}
\newcolumntype{R}[1]{>{\raggedleft\let\newline\\\arraybackslash\hspace{0pt}}m{#1}}



%Logo
%\AddToShipoutPictureBG{%
%  \AtPageUpperLeft{\raisebox{-\height}{\includegraphics[width=1.5cm]{uchicago.png}}}
%}

\newcolumntype{L}[1]{>{\raggedright\let\newline\\\arraybackslash\hspace{0pt}}m{#1}}
\newcolumntype{C}[1]{>{\centering\let\newline\\\arraybackslash\hspace{0pt}}m{#1}}
\newcolumntype{R}[1]{>{\raggedleft\let\newline\\\arraybackslash\hspace{0pt}}m{#1}}

\newcommand{\mr}{\multirow}
\newcommand{\mc}{\multicolumn}

%\newcommand{\comment}[1]{}

%Other parameters
\newcommand{\noutcomes}{95}
\newcommand{\noutcomesexpp}{357}
\newcommand{\noutcomesexpm}{343}
\newcommand{\noutcomesexpf}{355}
\newcommand{\treatsubsabc}{$75\%$}
\newcommand{\treatsubscarec}{$74\%$}
\newcommand{\treatsubscaref}{$63\%$}

%Counts
%Males
\newcommand{\positivem}{$78\%$}
\newcommand{\positivesm}{$29\%$}

%Females
\newcommand{\positivef}{$78\%$}
\newcommand{\positivesf}{$31\%$}

%Counts, control substitution
%Males
\newcommand{\positivecsnm}{$47\%$}
\newcommand{\positivescsnm}{$15\%$}

\newcommand{\positivecsam}{$79\%$}
\newcommand{\positivescsam}{$29\%$}

%Females
%% no alternative
\newcommand{\positivecsnf}{$84\%$}
\newcommand{\positivescsnf}{$55\%$}

%% alternative
\newcommand{\positivecsaf}{$79\%$}
\newcommand{\positivescsaf}{$33\%$}

%Pooled

%Effects
%Males

%Females
\newcommand{\empf}{$8$}
\newcommand{\yearsedf}{$1.7$}



%Pooled

%CBA
%IRR
%Males
\newcommand{\irrm}{$15\%$}
\newcommand{\irrsem}{$5\%$}

%Females
\newcommand{\irrf}{$9\%$}
\newcommand{\irrsef}{$7\%$}

%Pooled
\newcommand{\irrp}{$13\%$}
\newcommand{\irrsep}{$5\%$}

%BC
%Males
\newcommand{\bcm}{$11.24$}
\newcommand{\bcsem}{$4.60$}

%Females
\newcommand{\bcf}{$2.35$}
\newcommand{\bcsef}{$1.09$}

%Pooled
\newcommand{\bcp}{$5.63$}
\newcommand{\bcsep}{$2.15$}

%NPV streams
%Pooled
\newcommand{\parincomenpvp}{$\$119,346$}

\usepackage[stable]{footmisc}

\newcommand*\leftright[2]{%
  \leavevmode
  \rlap{#1}%
  \hspace{0.5\linewidth}%
  #2}

\newcommand{\orth}{\ensuremath{\perp\!\!\!\perp}}%
\newcommand{\indep}{\orth}%
\newcommand{\notorth}{\ensuremath{\perp\!\!\!\!\!\!\diagup\!\!\!\!\!\!\perp}}%
\newcommand{\notindep}{\notorth}

\externaldocument{abc_treatmenteffects_appendix_04-27-2017a_alz}
\pagenumbering{roman}

\begin{document}

\begin{titlepage}
\newgeometry{top=.8in, bottom=.8in, left=.8in, right=.8in}

\title{\Large \textbf{Gender Differences in the Benefits of a Prototypical Early Childhood Program: \\ Holding Tank}\thanks{This research was supported in part by: the Buffett Early Childhood Fund; the Pritzker Children's Initiative; the Robert Wood Johnson Foundation's Policies for Action program; the Leonard D. Schaeffer Center for Health Policy and Economics; NIH grants NICHD R37HD065072, NICHD R01HD054702, NIA R24AG048081, and NIA P30AG024968; Successful Pathways from School to Work, an initiative of the University of Chicago's Committee on Education and funded by the Hymen Milgrom Supporting Organization; the Human Capital and Economic Opportunity Global Working Group, an initiative of the Center for the Economics of Human Development and funded by the Institute for New Economic Thinking; and the American Bar Foundation. The views expressed in this paper are solely those of the authors and do not necessarily represent those of the funders or the official views of the National Institutes of Health. The authors wish to thank Frances Campbell, Craig and Sharon Ramey, Margaret Burchinal, Carrie Bynum, and the staff of the Frank Porter Graham Child Development Institute at the University of North Carolina Chapel Hill for the use of data and source materials from the Carolina Abecedarian Project and the Carolina Approach to Responsive Education. Years of partnership and collaboration have made this work possible. We thank Ruby Zhang for excellent research assistance. Collaboration with Andr\'{e}s Hojman, Yu Kyung Koh, Sylvi Kuperman, Stefano Mosso, Rodrigo Pinto, Joshua Shea, and Jake Torcasso on related work has strengthened the analysis in this paper. For helpful comments on various versions of the paper, we thank St\'{e}phane Bonhomme, Fl\'{a}vio Cunha, Steven Durlauf, David Figlio, Dana Goldman, Ganesh Karapakula, Sidharth Moktan, Tanya Rajan, Azeem Shaikh, Jeffrey Smith, Chris Taber, Matthew Tauzer, Ed Vytlacil, Jim Walker, and Matt Wiswall. We benefited from helpful comments received at the Leonard D. Schaeffer Center for Health Policy and Economics in December, 2016, and at the University of Wisconsin, February, 2017. For information on the implementation of the Carolina Abecedarian Project and the Carolina Approach to Responsive Education and assistance in data acquisition, we thank Peg Burchinal, Carrie Bynum, Frances Campbell, and Elizabeth Gunn. For information on childcare in North Carolina, we thank Richard Clifford and Sue Russell. The set of codes to replicate the computations in this paper are posted in a repository. Interested parties can request to download all the files. The address of the repository is \url{https://github.com/jorgelgarcia/abccare-cba}. To replicate the results in this paper, contact any of the authors, who will put you in contact with the appropriate individuals to obtain access to restricted data.}}

\author{
Jorge Luis Garc\'{i}a\\
Department of Economics\\
The University of Chicago \and
James J. Heckman \\
American Bar Foundation \\
Department of Economics\\
The University of Chicago \and
Anna L. Ziff \\
Center for the Economics of \\
Human Development \\
The University of Chicago}
\date{First Draft: January 5, 2016\\ This Draft: \today}

\maketitle
\restoregeometry
\end{titlepage}

\singlespacing
\begin{abstract}
\noindent This paper studies the life-cycle impacts of a widely emulated, high-quality intensive early childhood program with long-term follow up. The program starts early in life (at 8 weeks of age) and is evaluated by an RCT. There are multiple treatment effects which we summarize through interpretable aggregates. Girls have a greater number of statistically significant treatment effects than boys and effect sizes for them are generally bigger. The source of this difference is worse home environments for girls with greater scope for improvement by the program. Fathers of sons support their families more than fathers of daughters.

\end{abstract}

\noindent \textbf{Keywords}: Childcare, early childhood education, gender differences in responses to programs, health, long-term predictions, quality of life, randomized trials, substitution bias \\
\noindent \textbf{JEL codes}: J13, I28, C93\\


\bigskip
\begin{tabular}{ll}
Jorge Luis Garc\'{i}a                                       & James J. Heckman \\
Department of Economics                                & Department of Economics \\
University of Chicago                                       & University of Chicago \\
1126 East 59th Street                                     & 1126 East 59th Street \\
Chicago, IL 60637                                           & Chicago, IL 60637 \\
Phone: 773-677-7938                                     & Phone: 773-702-0634  \\
Email: jorgelgarcia@uchicago.edu                       & Email: jjh@uchicago.edu \\
                                                                       & \\
Anna L. Ziff                                         &  \\
Center for Economics of & \\
Human Development            & \\
University of Chicago                                        &  \\
1126 East 59th Street                       & \\
Chicago, IL 60637                                              &       \\
Phone: 773-834-0029                                    &  \\
Email: aziff@uchicago.edu                     &  \\

\end{tabular}

\clearpage

\restoregeometry
\doublespacing

\setcounter{page}{0}
\pagenumbering{arabic}



This seemingly contradictory result is explained by the types of treatment effects by gender and specific gender differences. For example, although the treatment reduced crime for both males and females, males tend to commit more costly crimes (e.g. violent crimes that have large costs to the victim).\footnote{\citet{Cohen-Bowles_2010_Estimating-Cost-Crime,Gregg_etal_2015_SocialRealities_BOOK}.} We explain the treatment effects in detail to help explain the gender differences present in these aggregate estimates.


\clearpage
\singlespacing
\bibliography{heckman}
\bibliographystyle{chicago}

\end{document}
